\documentclass[../../../include/open-logic-section]{subfiles}

\begin{document}

\olfileid{sth}{ord-arithmetic}{add}
\olsection{Panambahan Ordinal}

Upamané kita arep nambahaké $\alpha$ lan $\beta$. Kita bisa nyelehake
{salinan} $\beta$ persis sawisé salinan $\alpha$. (Kita perlu nganggo
\emph{salinan}, amarga saka
\olref[ordinals][basic]{ordinalsaresubsets} kita ngerti yèn $\alpha
\subseteq \beta$ utawa $\beta \subseteq \alpha$.) Kanthi intuisi, iki
tegesé ngliwati urutan langkah dawané $\alpha$, \emph{banjur}
ngliwati urutan langkah dawané $\beta$. Urutan kang diasilaké
bakal terurut becik; mula miturut
\olref[ordinals][ordtype]{thmOrdinalRepresentation} urutan iku
isomorfis karo ordinal (kang tunggal). Ordinal mau bisa dianggep
minangka \emph{gunggung} $\alpha$ lan $\beta$.

Iki gagasan intuitif ing waliking panambahan ordinal. Kanggo
netepaké kanthi ketat, kita miwiti saka gagasan nggawe
\emph{salinan} himpunan. Ing kéné kita nganggo tandha sawenang-wenang,
$0$ lan $1$, kanggo nandhani asal saben objek:

\begin{defn}\ollabel{defdissum}
\emph{Gunggung pisah} saka $A$ lan $B$ yaiku $A \disjointsum B = (A\times
\{0\}) \cup (B \times \{1\})$.
\end{defn}

Sabanjuré kita netepaké sawijining pengurutan tumrap pasangan ordinal:

\begin{defn}
Kanggo sembarang ordinal $\alpha_1, \alpha_2, \beta_1, \beta_2$, kandhaa yèn:
\begin{align*}
	\tuple{\alpha_1, \alpha_2} \rlexless \tuple{\beta_1, \beta_2}\text{ yen lan mung yen }&
	\text{utawa $\alpha_2 \in \beta_2$}\\
	& \text{utawa $\alpha_2 = \beta_2$ lan $\alpha_1 \in \beta_1$}
\end{align*}
\end{defn}

Iki pengurutan \emph{leksikografis kuwalik}, amarga unsur kapindho
diurutaké dhisik, banjur unsur kapisan. Élinga yèn kita arep
netepaké $\alpha \ordplus \beta$ minangka jinis urutan saka salinan
$\alpha$ kang diterusaké salinan $\beta$. Kanggo kuwi, kita kandha:

\begin{defn}\ollabel{defordplus}
Kanggo sembarang ordinal $\alpha$, $\beta$, gunggungé yaiku $\alpha \ordplus
\beta = \ordtype{\alpha \disjointsum \beta, \rlexless}$.
\end{defn}
\noindent
Wigatèkna yèn ing kéné kita rada nyimpang saka notasi kang ketat;
sakjané kita kudu nulis “$\Setabs{\tuple{x,y}\in \alpha\disjointsum\beta}{x \rlexless y}$” tinimbang “$\rlexless$”. Nanging, supaya ringkes,
ing pérangan sabanjuré kita tetep nganggo notasi iki.

Asil ing ngisor iki, bebarengan karo
\olref[ordinals][ordtype]{thmOrdinalRepresentation}, nandhesaké
yèn definisi kita wujudé trep:

\begin{lem}\ollabel{ordsumlessiswo}
$\tuple{\alpha \disjointsum \beta, \rlexless}$ iku pengurutan becik,
kanggo sembarang ordinal $\alpha$ lan $\beta$.
\end{lem}

\begin{proof}
Cetha yèn $\rlexless$ iku pengurutan kang nyambung ing $\alpha \disjointsum \beta$. Kanggo
nélakaké yèn saben subhimpunan ora kosong nduwèni unsur paling cilik, tetepna $X \subseteq \alpha
\disjointsum \beta$ kang ora kosong. Upamané $Y$ iku subhimpunan $X$ kang koordinat kapindhoné paling cilik, yaiku\ $Y = \Setabs{\tuple{\gamma, i} \in X}{(\forall \tuple{\delta, j} \in X)i \leq j}$. Saiki pilih unsur ing $Y$ kang koordinat kapisané paling cilik.
\end{proof}
\noindent
Dadi kita wis nduwèni definisi panambahan ordinal kang cetha
lan eksplisit. Iki sawijining kasunyatan kang ora nggumunaké
(élinga yèn $1 = \{0\}$, miturut
\olref[z][infinity-again]{defnomega}):
\begin{prop}
$\alpha \ordplus  1 = \ordsucc{\alpha}$, kanggo sembarang ordinal $\alpha$.
\end{prop}

\begin{proof}
Gatekna isomorfisme $f$ saka $\ordsucc{\alpha} = \alpha \cup
\{\alpha\}$ menyang $\alpha\disjointsum1 = (\alpha \times \{0\})
\cup (\{0\} \times \{1\})$ kang diwènèhaké déning $f(\gamma) =
\tuple{\gamma, 0}$ kanggo $\gamma \in \alpha$, lan $f(\alpha) = \tuple{0,
1}$.
\end{proof}
\noindent
Kajaba iku, gampang dituduhaké yèn panambahan manut sawetara
syarat rekursif:

\begin{lem}\ollabel{ordadditionrecursion}
Kanggo sembarang ordinal $\alpha, \beta$, kita nduwèni:
\begin{align*}
	\alpha\ordplus 0 &= \alpha\\
	\alpha \ordplus  (\beta\ordplus 1) &= (\alpha \ordplus  \beta) \ordplus  1\\
	\alpha  \ordplus  \beta &= \supstrict_{\delta < \beta}(\alpha \ordplus  \delta) && \text{yèn $\beta $ ordinal limit}
\end{align*}
\end{lem}

\begin{proof}
Kita priksa saben kasus; kapisan:
\begin{align*}
	\alpha \ordplus  0
	& = \ordtype{(\alpha \times \{0\}) \cup (0 \times \{1\}), \rlexless} \\
	&= \ordtype{(\alpha \times \{0\}) \cup \emptyset, \rlexless}\\
	&= \alpha\\
	\alpha \ordplus (\beta \ordplus  1)
	%&= \ordtype{(\alpha\times \{0\}) \cup (\ordtype{\beta \ordplus  1}\times \{1\}), \rlexless} \\
	&= \ordtype{(\alpha\times \{0\}) \cup (\ordsucc{\beta}\times \{1\}), \rlexless} \\
	&= \ordtype{(\alpha\times \{0\}) \cup (\beta \times \{1\}), \rlexless} \ordplus 1\\
	&= (\alpha \ordplus  \beta) \ordplus  1
\end{align*}
Saiki upamané $\beta \neq \emptyset$ iku ordinal limit. Yèn $\delta < \beta$, uga
$\delta\ordplus 1 < \beta$, mula $\alpha \ordplus \delta$ iku
segmen awal sejati saka $\alpha \ordplus \beta$. Mulané $\alpha
\ordplus \beta$ iku wates ndhuwur ketat kanggo $X = \Setabs{\alpha
\ordplus \delta}{\delta < \beta}$. Kajaba iku, yèn $\alpha \leq \gamma <
\alpha \ordplus \beta$, mesthi $\gamma = \alpha \ordplus
\delta$ kanggo sawatara $\delta < \beta$. Mula $\alpha \ordplus \beta =
\supstrict_{\delta< \beta}(\alpha\ordplus \delta)$.
\end{proof}

Nanging ana kasunyatan kang narik kawigaten. Kanggo netepaké
panambahan ordinal, kita bisa \emph{kosokbaliné} mung nganggo
Teorema Rekursi Transfinit, lan nyatakaké persamaan rekursi
persis kaya ing \olref{ordadditionrecursion} (nanging nganggo
“$\ordsucc{\beta}$” tinimbang “$\beta \ordplus 1$”).

Dadi ana rong cara béda kanggo netepaké operasi ing ordinal.
Kita bisa netepaké kanthi \emph{sintetis}, yaiku mbangun himpunan
terurut becik kanthi eksplisit lan njupuk jinis urutané. Utawa
kita bisa netepaké kanthi \emph{rekursif}, mung kanthi nyatakaké
persamaan rekursiné. Yèn ditindakake kanthi bener, asile padha.
Awit \olref[ordinals][ordtype]{thmOrdinalRepresentation} njamin
yèn persamaan rekursi iki nemtokaké ordinal kang \emph{tunggal}.

Ing akèh babagan, aritmetika ordinal tumindaké padha karo
panambahan wilangan asli. Contoné, kita bisa mbuktèkaké iki:

\begin{lem}\ollabel{ordinaladditionisnice}
Yèn $\alpha, \beta, \gamma$ ordinal, mula:
\begin{enumerate}
	\item\ollabel{ordaddition1} yèn $\beta < \gamma$, mula $\alpha
	\ordplus  \beta < \alpha \ordplus  \gamma$
	\item\ollabel{ordaddition2} yèn $\alpha \ordplus  \beta =
	\alpha\ordplus \gamma$, mula $\beta = \gamma$
	\item\ollabel{ordaddition3} $\alpha \ordplus  (\beta \ordplus
	\gamma) = (\alpha \ordplus  \beta) \ordplus  \gamma$, yaiku
	panambahan asosiatif
	\item\ollabel{ordaddition4} yèn $\alpha \leq \beta$, mula $\alpha
	\ordplus  \gamma \leq \beta \ordplus \gamma$
\end{enumerate}
\end{lem}

\begin{proof}
Kita mbuktèkaké \olref{ordaddition3}; sisane dadi latihan. Buktiné
nganggo Induksi Transfinit Prasaja marang $\gamma$, kanthi
\olref{ordadditionrecursion}. Nalika $\gamma = 0$:
\[
(\alpha \ordplus  \beta) \ordplus  0 = \alpha \ordplus  \beta  = \alpha \ordplus  (\beta \ordplus  0)
\]
Nalika $\gamma = \delta\ordplus 1$, upamané kanggo induksi yèn $(\alpha
\ordplus  \beta) \ordplus  \delta = \alpha \ordplus  (\beta \ordplus
\delta)$; saiki nganggo \olref{ordadditionrecursion} ping telu:
\begin{align*}
	(\alpha \ordplus  \beta) \ordplus  (\delta \ordplus  1) & = ((\alpha \ordplus  \beta) \ordplus  \delta)\ordplus 1\\
	& = (\alpha \ordplus  (\beta \ordplus  \delta)) \ordplus  1\\
	& = \alpha \ordplus  ((\beta \ordplus  \delta)\ordplus 1)\\
	& = \alpha \ordplus  (\beta \ordplus  (\delta\ordplus 1))
\end{align*}
Nalika $\gamma$ iku ordinal limit, upamané kanggo induksi yèn
$\delta \in \gamma$ ngakibataké $(\alpha \ordplus  \beta) \ordplus  \delta =
\alpha \ordplus  (\beta \ordplus  \delta)$; mula:
\begin{align*}
	(\alpha \ordplus  \beta) \ordplus  \gamma & = \supstrict_{\delta < \gamma}((\alpha \ordplus  \beta) \ordplus  \delta) \\
	&= \supstrict_{\delta < \gamma}(\alpha \ordplus  (\beta \ordplus  \delta))\\
	&= \alpha \ordplus  \supstrict_{\delta < \gamma}(\beta \ordplus  \delta)\\
	& = \alpha \ordplus  (\beta \ordplus  \gamma)
\end{align*}
\end{proof}

\begin{prob}
Buktèkna pérangan liya saka
\olref[sth][ord-arithmetic][add]{ordinaladditionisnice}.
\end{prob}

Miturut perkara-perkara iki, panambahan ordinal mesthiné wis
kulina tumrap kita. Nanging ana siji bab wigati kang ndadèkaké
panambahan ordinal \emph{ora} kaya panambahan wilangan asli.

\begin{prop}\ollabel{ordsumnotcommute}
Panambahan ordinal {ora} komutatif; $1 \ordplus  \omega = \omega <
\omega \ordplus  1$.
\end{prop}

\begin{proof}
Wigatèkna yèn $1 \ordplus  \omega = \supstrict_{n < \omega} (1 \ordplus n)
= \omega \in \omega \cup \{\omega\} = \ordsucc{\omega} = \omega
\ordplus  1$.
\end{proof}
\noindent
Sanadyan wiwitané iki bisa nggumunaké, \emph{sakjané ora perlu}.
Nalika mikiraké $1 \ordplus \omega$, kita mikiraké jinis urutan kang
kapikolehi kanthi nyelehake unsur tambahan \emph{sadurungé}
kabeh wilangan asli. Kaya panalaran babagan Hotel Hilbert ing
\olref[sfr][infinite][hilbert]{sec}, kanthi intuisi unsur tambahan
ing wiwitan iki ora ngowahi jinis urutan sakabèhé. Nanging nalika
mikiraké $\omega \ordplus 1$, kita mikiraké jinis urutan kang
kapikolehi kanthi nyelehake unsur tambahan \emph{sawisé}
kabeh wilangan asli. Asilé béda banget!

\end{document}
